\chapter{Introduction}

\section{Clinical Motivation}

\subsection{Why CT-Based Mucus Plug Assessment Matters}
Mucus plugs are intraluminal accumulations of abnormally retained mucus that partially or completely obstruct the airways \cite{fahy2010airwaymucus,dunican2018autopsy}. In healthy lungs, mucus contributes to airway defense and is normally cleared through coordinated mucociliary transport \cite{fahy2010airwaymucus}. In chronic airway disease altered mucus composition, impaired clearance, and inflammatory processes can transform a normally protective material into a persistent source of obstruction \cite{fahy2010airwaymucus,dunican2018autopsy}. In asthma, autopsy and imaging studies indicate that mucus plugging is not restricted to rare terminal events, but can also be present as a clinically relevant feature of chronic disease \cite{dunican2018autopsy}. Greater CT-measured mucus plug burden has been associated with worse airflow obstruction and with eosinophilic or type 2 airway inflammation, which makes mucus plugging relevant not only as an imaging finding but also as an indicator of disease burden \cite{dunican2018mucus}.


Computed tomography (CT) is relevant in this context because it provides noninvasive visualization of airway and lung structure and has long been used to study structural changes related to respiratory dysfunction \cite{dejong2005ctairways}. This is important for mucus plugging because the abnormality is distributed within the bronchial tree rather than confined to a single projection image or laboratory measurement. On CT, mucus plugs may appear as complete bronchial occlusions seen as tubular, branching, or rounded opacities that can be followed across sequential slices and related to the bronchial lumen \cite{dunican2018autopsy,dunican2018mucus}. CT therefore supports structured assessment of plugging throughout a full volumetric examination, rather than limiting interpretation to isolated image examples or local observations \cite{dejong2005ctairways,dunican2018mucus}. Raw axial CT slices from the study data are shown in Figure~\ref{fig:raw_ct_examples} to illustrate the imaging context in which mucus plug assessment is performed.

\begin{figure}[H]
    \centering
    \includegraphics[width=0.32\textwidth]{figures/rawct1.png}\hfill
    \includegraphics[width=0.32\textwidth]{figures/rawct2.png}\hfill
    \includegraphics[width=0.32\textwidth]{figures/rawct3.png}
    \caption[Representative raw chest CT slices from the study data]{Representative raw axial chest CT slices from the study data shown in lung windowing. The figure illustrates the imaging context in which mucus plug assessment is performed and emphasizes that clinically relevant information is embedded within a volumetric examination rather than a single isolated image.}
    \label{fig:raw_ct_examples}
\end{figure}

A patient-level mucus score is therefore meaningful as a compact summary of disease-relevant information across the examination. In a segment-based CT scoring approach, the lungs are assessed systematically and the presence or absence of mucus plugging is summarized into an aggregate score over the full scan \cite{dunican2018mucus}. Such a score is clinically useful because it reflects overall burden at the level at which patients are followed, compared, and interpreted in relation to lung function and inflammatory status \cite{dunican2018mucus}. At the same time, deriving this score manually requires systematic review across multiple bronchopulmonary segments and sequential CT slices, which makes the process time-consuming and dependent on careful volumetric interpretation.

The clinical motivation for this thesis follows from this setting. Existing CT-based mucus plug work shows that the target is clinically relevant, but leaves open how it can be estimated computationally under realistic data constraints. The purpose of this thesis is therefore not to replace clinical judgment with a single image-level shortcut, but to investigate whether chest CT can support patient-level scoring through a constrained and reproducible machine learning formulation.

\subsection{Problem Statement}

This thesis addresses automated patient-level prediction of mucus plug burden from chest CT. The central technical difficulty is that the clinically meaningful target is defined for the patient as a whole, whereas the available image evidence is observed at slice level. The task is therefore not to classify one image in isolation, but to infer a single patient-level score from a volumetric examination in which relevant findings may appear unevenly across slices.

This creates a weakly supervised problem setting. Weak supervision arises when the available labels are coarser than the image evidence used for learning \cite{ren2023weakly}. In the present case, the model processes slice-level image content, but the available supervision is defined only at the patient level. The task is therefore more challenging than settings in which the training labels directly identify the relevant image content, because the model must learn how local observations contribute to a patient-level target without slice-level annotations \cite{ren2023weakly,qi2021drmil}.

The problem is further constrained by data availability. In a small internal cohort, not every slice is equally informative for patient-level prediction, and model design choices become closely tied to the risk of overfitting and unstable generalization \cite{dhiman2023samplesize}. Existing clinical CT scoring literature establishes that mucus plug burden can be assessed meaningfully at patient level, while related weakly supervised CT literature provides methodological ideas for learning from coarse case-level labels. What remains less well addressed is the combination of these perspectives, a patient-level mucus score prediction from chest CT under weak supervision and limited data. This thesis is positioned within that gap.


\section{Thesis Aim and Research Questions}

\subsection{Objective}
The objective of this thesis is to investigate whether patient-level mucus plug burden can be estimated from chest CT using a reproducible machine learning pipeline that remains realistic with respect to the available supervision and dataset size. The thesis studies how representative slice selection, CT input representation, feature fusion, compact CNN backbones, and transfer-oriented evaluation can support patient-level prediction from volumetric CT. The transfer experiments use MosMed as an external CT dataset for representation transfer and domain-shift analysis, not as external validation of mucus-burden prediction, because MosMed has a different COVID-related severity target.

The aim is therefore not to propose a new clinical scoring system, nor to claim a clinically validated decision-support tool. Instead, the thesis seeks to evaluate a constrained computational approach to a clinically meaningful CT scoring problem and to clarify which design choices are useful, defensible, and informative within this setting.


\subsection{Research Questions}
The thesis is guided by the following research questions:

\begin{description}
    \item[\textbf{RQ1.}] To what extent can a slice-based weakly supervised regression pipeline predict patient-level mucus plug burden from chest CT?
    
    \item[\textbf{RQ2.}] How does representative slice selection affect patient-level prediction compared with simpler slice sampling?
    
    \item[\textbf{RQ3.}] How do alternative CT input representations and feature-fusion strategies, including CNN-based learned representations and bag-of-visual-words-style handcrafted features, compare under limited-data conditions?
    
    \item[\textbf{RQ4.}] How do compact pretrained CNN backbones compare for patient-level mucus plug burden prediction from chest CT?
    
    \item[\textbf{RQ5.}] What do transfer-learning experiments with external chest CT data and a different target reveal about domain shift, adaptation, and the limits of generalization in this small-cohort setting?
\end{description}

These research questions are formulated to be empirically answerable within the scope of the thesis. They connect the clinical motivation for patient-level mucus scoring to the methodological decisions evaluated in the later chapters.


\subsection{Scope and Delimitations}
The scope of the thesis is limited to patient-level prediction from chest CT. It does not perform segmentation of mucus plugs, localization of individual plugged airways, or detection of specific bronchopulmonary segments. It also does not include full 3D model training, large-scale training, or a clinically ready decision-support system. The contribution is a reproducible machine learning pipeline and empirical analysis of its design choices, not a validated clinical product.

The main delimitations follow directly from the data and project scope. The available labels are patient-level scores, so the thesis cannot make slice-level claims about where mucus plugs are located unless such information is introduced in future annotation work. The study is also conducted under limited-data conditions, which necessarily constrains the strength of claims about generalization beyond the present setting \cite{dhiman2023samplesize}. The experiments evaluate design choices within the implemented pipeline, they do not exhaust the full space of possible medical-imaging models or clinical predictors.


\section{Contribution Summary}

Given these boundaries, the thesis makes the following contributions:

\begin{itemize}
    \item it formulates patient-level mucus plug burden prediction from chest CT as a weakly supervised regression problem.
    
    \item it develops a reproducible slice-based prediction pipeline that separates preprocessing, representative slice selection, input representation, feature extraction, and patient-level aggregation under limited-data conditions.
    
    \item it evaluates how key design choices, including slice-selection strategy, CT input representation, feature fusion, CNN backbone choice, aggregation, and transfer-oriented evaluation, affect patient-level prediction performance.
    
    \item it provides an empirical analysis of the strengths and limitations of this formulation in relation to weak supervision, small-cohort learning, and clinically relevant patient-level scoring.
\end{itemize}

The contribution of the thesis is therefore methodological and empirical rather than clinical in the narrow sense. It investigates a comparatively underexplored combination of target, supervision setting, and computational design within a reproducible machine learning framework.


\section{Thesis Outline}
The remainder of the thesis is organized as follows. Chapter~2 reviews the literature most relevant to this work, with emphasis on CT-based mucus plug assessment, learning approaches for patient-level CT analysis, and the research gap addressed by the thesis. Chapter~3 presents the theory and problem formulation, including the patient-level prediction task, the label structure, weak supervision, slice representation, feature learning, aggregation, regression, and evaluation. Chapter~4 describes the implemented methodology, covering dataset and target definition, preprocessing, representative slice selection, input representation, model architecture, patient-level aggregation, training, transfer learning, and reproducibility. Chapter~5 presents the experimental protocol and quantitative results from the internal cross-validation and transfer-learning experiments. Chapter~6 discusses the findings in relation to the research questions, including the ablation-style interpretation of pipeline components, weak supervision, low and negative \(R^2\) values, transfer/domain-shift results, and the main limitations of the study. Finally, Chapter~7 concludes the thesis by summarizing the main contributions, limitations, future directions, and final remarks.


